\documentclass[11pt,a4paper]{article}
\usepackage[margin=2.5cm]{geometry}
\usepackage[T1]{fontenc}
\usepackage{amsmath,amssymb,booktabs,graphicx,tabularx,array}
\usepackage[authoryear,round]{natbib}
\usepackage[colorlinks=true,linkcolor=blue,citecolor=blue,urlcolor=blue]{hyperref}
\graphicspath{{results/}}
\input{results/preflight_numbers.tex}
\emergencystretch=3em
\title{Ecological assets behind improving accounts:\\a boundary-matched test of composition masking}
\author{Working manuscript prepared for Amin Abaee\\\small F01: staged empirical implementation}
\date{29 October 2026}
\begin{document}
\maketitle
\begin{center}\fbox{\parbox{0.92\textwidth}{\small\textbf{Stage 1 working draft, not a completed empirical paper.} The source audit and tested computational kernels reported below have been implemented. Physical-stock trends, conversion associations, diagnostic performance and divergence prevalence have not yet been estimated. Author verification and approval are required.}}
\end{center}
\begin{abstract}
Biological-supply accounts and physical ecological assets describe different quantities. We develop a forest-led empirical study asking whether changes in national biocapacity accounts agree with independently estimated forest extent and above-ground woody biomass on matched boundaries and intervals. The design retains concordant, divergent and unresolved outcomes, distinguishes conversion from other disturbances, and evaluates spatially offsetting local net changes without treating them as gross biological flows. This implementation draft reports the first data-access and eligibility stage, not ecological trend results. A frozen national workbook contains \NationalRows\ country-year records across \CountryIDs\ identifiers; \AccountPass\ identifiers pass the account-only 2015--2023 preflight. All eight deliberately selected pilot candidates pass that limited screen, but their physical support and validation gates remain pending. Small version-frozen biomass files confirm access and expose consequential interpretation issues: nominal spatial resolution is not equal-area support, aggregated density is per grid-cell rather than forest area, and change-quality labels derived from the same endpoint maps are not independent validation truth. Tested area, decomposition and cancellation kernels establish computational identities but do not establish measurement validity. Subsequent stages must resolve boundaries, coverage, change uncertainty, calibration and transition validation before an empirical account--asset conclusion is warranted.
\end{abstract}

\section{Introduction}
A biocapacity account normalizes biological production-related quantities into global hectares. A physical biomass stock, in contrast, measures a specified mass over a defined support. Neither is an automatic measure of the other. The distinction between stocks and flows in Ecological Footprint interpretation has already been developed in prior work \citep{Mancini2017}; the present study does not claim to discover it.

The theoretical ECOMOD companion demonstrates conditional mechanisms through which conversion to a more annually productive land class can raise a model account while capital-land area falls. Its C13 revision quantifies static identification limits and within-model contribution budgets. Those calculations do not identify physical assets in national accounts or establish an empirical conversion mechanism. A follow-up therefore needs independently estimated physical changes and independently validated land transitions, rather than another algebraic residual or uncalibrated ecological index.

Spatial compensation is also not a new concept. European forest-condition accounting provides a closely related example in which an improving aggregate can coexist with declining condition in part of the domain \citep{Maes2023}. Forest condition, biomass, extent and regeneration remain distinct outcomes. The prospective contribution here is a reproducible, boundary-matched comparison of account signals and specified asset changes, with explicit uncertainty, coverage and incremental-information tests. It is not a prior claim of global decline, ecological safety, causality, or first-ever novelty.

\section{Questions and staged design}
The primary question is whether national account changes and physical forest changes agree on the same country--interval support. Secondary questions concern spatially offsetting local net changes, associations with independently observed forest-to-cropland conversion, and whether additional reporting information improves a declared diagnostic at comparable error and coverage. The sample will not be restricted to increasing-account countries.

The primary candidate interval is 2015--2023. A 2010--2020 sensitivity is conditional on the documented cross-sensor change limitations. The account vintage is NEFBA Edition 2026, results and metadata v1.0 \citep{Lo2026}. The physical source under examination is ESA CCI Biomass v7.0 \citep{SantoroCartus2026}. Earlier and later product releases will not be mixed to manufacture an observational time series.

The initial pilot candidates are Brazil, Indonesia, Democratic Republic of the Congo, Cameroon, Canada, Finland, France and Australia. They represent deliberate logistical and environmental variety, not a probability sample or an outcome-selected set of masking cases. Territorial definitions, overseas areas, source masks, valid observation support, forest definitions, and availability of independent validation remain gates. The accessible account population is not automatically the final empirical study population.

\begin{table}[htbp]\centering\small
\caption{Implementation stages and the evidence needed before progressing.}
\begin{tabularx}{\textwidth}{p{0.17\textwidth}X}
\toprule
Stage & Deliverable and gate\\\midrule
1: preflight & Frozen sources, account availability, coordinate/schema audit, tested kernels, working protocol and manuscript. Completed within the limited scope reported here.\\
2: physical pilot & Boundary-matched support, correct density denominator, missing/converted-pixel accounting, source lineage and independently justified change precision. Pending.\\
3: calibration & Noisy-null and misspecification tests, tuning/evaluation separation, independent references and identical-support reporting-package contrasts. Pending.\\
4: estimation & Uncertainty-aware stock/extent/density and spatial-net-change results; conversion associations only with validated transitions. Pending.\\
5: completion & Results-driven main text/SI, reference checks, independent rerun and author approval. Pending.\\\bottomrule
\end{tabularx}\end{table}

Soil, wetland and assessed-fish extensions are conditional modules, not mandatory rows in a superficially complete global annual panel. No scalar ecological-capital total, depletion horizon, physiological interpretation of a remote-sensing index, or calibration of ECOMOD's quality variable is part of the default design.

\section{Observation operators and estimands}
\subsection{Accounts, area and biomass}
The primary account is absolute total biocapacity. A terrestrial subtotal, excluding-built-up variants, the forest component, and per-person measures will be reported separately. National account categories are not forced into the conserved two-land partition of the theoretical model. The workbook's own category areas are useful account inputs but are not independent physical validation of that account.

For a defined geographic support $\Omega$, the stock calculation is
\[
Q_t=\int_{\Omega} d_t(x)\,\mathrm d A(x),
\]
where the density denominator and observed/missing support must be specified. The downloaded coarse CCI aggregate is an average per grid-cell area, not per forest area, as clarified in its Product User Guide \citep{PUG2026}. Multiplying that quantity by forest area again would apply an unjustified additional mask. Forest-only stocks and densities require an independent support definition and a justified mapping. A globally rectangular coordinate array does not establish globally valid observation coverage.

Cell areas must follow the raster bounds. The tested kernel integrates rectangles bounded by meridians and parallels on the WGS84 ellipsoid, accounts for latitude dependence, and rejects unsplit dateline crossings. It does not treat nominal 50 km geographic pixels as identical squares.

\subsection{Extent--density decomposition}
For positive forest areas $A_0,A_1$ and compatible stocks $Q_0,Q_1$, define $D_t=Q_t/A_t$. The exact symmetric decomposition is
\begin{align*}
\Delta Q_{\mathrm{extent}}&=(A_1-A_0)(D_1+D_0)/2,\\
\Delta Q_{\mathrm{density}}&=(D_1-D_0)(A_1+A_0)/2.
\end{align*}
Their sum is $Q_1-Q_0$. These are descriptive contributions, not causal effects. The supplied function refuses zero-area endpoints rather than assigning an unobserved counterfactual density. Appearing or disappearing strata require a separate convention and sensitivity analysis.

\subsection{Spatial cancellation and temporal evidence}
Given comparable local stock changes $\Delta Q_i$, let
\[
G=\sum_i[\Delta Q_i]_+,\quad L=\sum_i[-\Delta Q_i]_+,\quad
K=1-\frac{|G-L|}{G+L}.
\]
The offset mass is $\min(G,L)$; $K$ is undefined when $G+L=0$. $G$ and $L$ describe positive and negative \emph{local net changes}, not observed gross growth, mortality, harvest or recruitment. Spatial aggregation changes them; noisy zero-change estimates can produce large apparent cancellation. Effective support and noisy-null calibration are therefore prerequisites, not optional cosmetic sensitivities. Temporal paths will be compared with endpoints and simple trend/disturbance baselines only where actual repeat observations support the comparison.

\section{Uncertainty, lineage and reporting comparisons}
Change uncertainty must include temporal and spatial dependence, systematic bias, land-transition error, and missing/converted-pixel treatment. A set of per-pixel standard deviations cannot be summed as independent errors without justification. Validation of stock levels is not automatically validation of stock change. Independent transition validation requires a defensible reference design; mapped tree-cover loss alone is not mapped forest-to-cropland conversion \citep{Olofsson2014}.

Three operational restrictions implement the earlier reviews. First, inputs are audited for availability, target derivation and shared training/reference information. Retrospective observation is not automatically a real-time forecast task, but target-derived quality flags cannot serve as independent truth. Second, reporting packages must be rescored on the same cases and weights, with pairwise overlap and assessed coverage also disclosed. Third, proposed added features are checked for algebraic redundancy and physical identification; improved prediction alone is not proof of a new state or mechanism.

Candidate rules and gate ablations require a whole-pipeline simulation/validation stage with separate tuning and evaluation and Monte Carlo uncertainty. Provider quality labels are not automatically a 95\% test. Practical-change thresholds and final classification rules remain to be specified in a documented amendment after blinded calibration; no ecological safety threshold is imported from the account balance or from a provider code constant.

\section{Implemented Stage 1 results}
\subsection{National account preflight}
The retrieved workbook contains \NationalRows\ national records, 51 columns and \CountryIDs\ country/territory identifiers spanning 1961--2025. The window audit requires every year from 2015 through 2023, appropriate publication flags and finite, nonnegative required values. Flags \texttt{+++}, \texttt{++} and \texttt{+} all publish the biocapacity quantities needed here; their differing trade/consumption detail is not graded ecological accuracy.

Of the \CountryIDs\ identifiers, \AccountPass\ pass this account-only preflight and \AccountFail\ do not (Figure~\ref{fig:coverage}). All eight pilot candidates pass the limited account screen. These counts do not assert physical comparability, an eligible global sample or ecological trends. The largest five-component account-total discrepancy within the primary window is 2 gha, within the predeclared maximum 3 gha rounding tolerance. No paired country account-change signs were computed.

\begin{figure}[htbp]\centering
\includegraphics[width=\textwidth]{account_coverage.pdf}
\caption{Account-only availability by source UN region. Passing the screen is not physical-data eligibility, independent validation, or a trend result.}\label{fig:coverage}
\end{figure}

\subsection{Biomass access and consequential interpretation checks}
Two small v7 files were successfully retrieved and hashed: an 18-epoch nominal 50 km biomass aggregate and a 2010--2020 change example. Only headers, coordinates, time values and metadata were read; biomass/change pixel values and signs were not inspected. The files are access/schema smoke tests, not national stock estimates.

The aggregate has 720 longitude by 360 latitude cells at 0.5\textdegree\ spacing and includes biomass and uncertainty layers. It must not be treated as equal-area sampling. The change file references an absent \texttt{time\_bnds1} variable and does not encode its quality-flag meanings internally. These issues require explicit ingestion/documentation handling, not silent repair or rejection of all source data.

Targeted reading of the v7 guide and inspection---not execution---of the provider change-generation script supplied additional constraints. The guide cautions against uncritical interpretation of map differences, describes density per total grid-cell area and documents a spatial-covariance contribution to aggregate uncertainty. The example code's change-SD calculation omits a temporal covariance term, and its quality rules use the same endpoint maps and their uncertainties plus further rule/filter operations. These flags are therefore target-derived diagnostics, not independent reference observations or a calibrated 95\% decision rule. The final pilot must examine their fitness rather than copy their labels into a ground-truth column.

\subsection{Computational checks}
Thirteen unit tests passed. They include WGS84 total-area and partition checks, latitude dependence, pure-extent and pure-density recovery, 500 synthetic mixed decompositions, invalid/zero-area rejection, and cancellation edge cases. A deterministic noisy-null smoke test confirms that large point-estimate cancellation can arise without true asset change. It is not an operating-characteristic calibration experiment. Unit identities establish computational consistency, not the truth of ecological measurements.

\section{Current inference limit and next implementation}
Stage 1 establishes account-data access, a reproducible availability screen and several necessary interpretation safeguards. It establishes neither account--asset divergence nor its absence. The next release must complete the territorial/support and change-validation pilot, then calibrate and run the empirical estimators. If physical change is unresolved, the paper must retain abstention or narrow its claims rather than manufacture a positive mechanism finding.

The historical planning and world-level account audits were known in advance. During source discovery, headers and earliest Afghan historical rows were inspected; the present specification is an openly disclosed working protocol, not an externally registered or pristine prospective preregistration. Later changes require recorded amendments. This manuscript will be revised from actual pilot/analysis tables, not filled with anticipated results.

\section*{Availability, attribution and AI assistance}
The F01 package supplies the working configuration/protocol, source hashes and retrieval instructions, derived availability tables, schema audit, tested kernels, source-to-result pipeline and this manuscript. Raw CCI files and restricted document contents are not redistributed. NEFBA-derived availability/metadata outputs carry the workbook's CC BY-SA 4.0 attribution and share-alike requirements; ESA CCI and the Biomass CCI project and providers are acknowledged under their data terms. No claim of journal submission, complete novelty clearance or author approval is made. Arena.ai Agent Mode assisted with this working protocol, source inspection, code, tests and draft writing; author review is required.

\begin{thebibliography}{9}
\bibitem[Mancini et al.(2017)]{Mancini2017} Mancini, M. S., et al. (2017). Stocks and flows of natural capital: Implications for Ecological Footprint. \emph{Ecological Indicators}. \url{https://doi.org/10.1016/j.ecolind.2017.01.033}.
\bibitem[Maes et al.(2023)]{Maes2023} Maes, J., et al. (2023). Accounting for forest condition in Europe based on an international statistical standard. \emph{Nature Communications}. \url{https://doi.org/10.1038/s41467-023-39434-0}.
\bibitem[Olofsson et al.(2014)]{Olofsson2014} Olofsson, P., et al. (2014). Good practices for estimating area and assessing accuracy of land change. \emph{Remote Sensing of Environment}. \url{https://doi.org/10.1016/j.rse.2014.02.015}.
\bibitem[Lo et al.(2026)]{Lo2026} Lo, K., et al. (2026). \emph{National Ecological Footprint and Biocapacity Accounts, Edition 2026 (1961--2025): Results and metadata, v1.0}. Prepared for Footprint Data Foundation by researchers at York University and the University of Iceland. \url{https://footprint.info.yorku.ca/data/}.
\bibitem[Santoro and Cartus(2026)]{SantoroCartus2026} Santoro, M., and Cartus, O. (2026). ESA Biomass Climate Change Initiative: Global datasets of forest above-ground biomass for 2005--2012 and 2015--2024, v7.0. CEDA. \url{https://doi.org/10.5285/6429d1aafe1e43b9b414e4a5a7f8b903}.
\bibitem[CCI Biomass(2026)]{PUG2026} CCI Biomass (2026). \emph{Product User Guide v7.0}, 19 March 2026. Targeted sections 2.3, 4--6. \url{https://climate.esa.int/en/projects/biomass/key-documents/}. Associated provider change script inspected, not executed; its exact hash is supplied in the source-method register.
\end{thebibliography}
\end{document}
